% OpenStax Calculus Volume 1 -- Section 4.10 Exercises: Antiderivatives
% Exercises reproduced from OpenStax, licensed under CC BY 4.0.
% Source: https://openstax.org/books/calculus-volume-1/pages/4-10-antiderivatives
\documentclass{exercises}
\title{OpenStax Calculus Volume 1}
\subtitle{Section 4.10 Exercises: Antiderivatives}
\sourceurl{https://openstax.org/books/calculus-volume-1/pages/4-10-antiderivatives}
\author{}
\date{}

\begin{document}
\maketitle


For the following exercises, show that $F(x)$ are antiderivatives of $f(x)$.

\begin{enumerate}[start=1]
\item $F(x)=5x^{3}+2x^{2}+3x+1,f(x)=15x^{2}+4x+3$
\item $F(x)=x^{2}+4x+1,f(x)=2x+4$
\item $F(x)=x^{2}e^{x},f(x)=e^{x}(x^{2}+2x)$
\item $F(x)=\cos x,f(x)=-\sin x$
\item $F(x)=e^{x},f(x)=e^{x}$
\end{enumerate}

For the following exercises, find the general form for the antiderivative of the function.

\begin{enumerate}[resume]
\item $f(x)=\frac{1}{x^{2}}+x$
\item $f(x)=e^{x}-3x^{2}+\sin x$
\item $f(x)=e^{x}+3x-x^{2}$
\item $f(x)=x-1+4\sin(2x)$
\end{enumerate}

For the following exercises, find the general form for the antiderivative $F(x)$ of each function $f(x)$.

\begin{enumerate}[resume]
\item $f(x)=5x^{4}+4x^{5}$
\item $f(x)=x+12x^{2}$
\item $f(x)=\frac{1}{\sqrt{x}}$
\item $f(x)={(\sqrt{x})}^{3}$
\item $f(x)=x^{1/3}+{(2x)}^{1/3}$
\item $f(x)=\frac{x^{1/3}}{x^{2/3}}$
\item $f(x)=2\sin(x)+\sin(2x)$
\item $f(x)=\sec^{2}(x)+1$
\item $f(x)=\sin x\cos x$
\item $f(x)=\sin^{2}(x)\cos(x)$
\item $f(x)=0$
\item $f(x)=\frac{1}{2}\csc^{2}(x)+\frac{1}{x^{2}}$
\item $f(x)=\csc x\cot x+3x$
\item $f(x)=4\csc x\cot x-\sec x\tan x$
\item $f(x)=8\sec x(\sec x-4\tan x)$
\item $f(x)=\frac{1}{2}e^{-4x}+\sin x$
\end{enumerate}

For the following exercises, evaluate the integral.

\begin{enumerate}[resume]
\item $\int(-1)\,dx$
\item $\int \sin x\,dx$
\item $\int(4x+\sqrt{x})\,dx$
\item $\int \frac{3x^{2}+2}{x^{2}}\,dx$
\item $\int(\sec x\tan x+4x)\,dx$
\item $\int(4\sqrt{x}+\sqrt[4]{x})\,dx$
\item $\int(x^{-1/3}-x^{2/3})\,dx$
\item $\int \frac{14x^{3}+2x+1}{x^{3}}\,dx$
\item $\int(e^{x}+e^{-x})\,dx$
\end{enumerate}

For the following exercises, solve the initial value problem.

\begin{enumerate}[resume]
\item $f'(x)=x^{-3},f(1)=1$
\item $f'(x)=\sqrt{x}+x^{2},f(0)=2$
\item $f'(x)=\cos x+\sec^{2}(x),f(\frac{\pi}{4})=2+\frac{\sqrt{2}}{2}$
\item $f'(x)=x^{3}-8x^{2}+16x+1,f(0)=0$
\item $f'(x)=\frac{2}{x^{2}}-\frac{x^{2}}{2},f(1)=0$
\end{enumerate}

For the following exercises, find two possible functions $f$ given the second- or third-order derivatives.

\begin{enumerate}[resume]
\item $f''(x)=x^{2}+2$
\item $f''(x)=e^{-x}$
\item $f''(x)=1+x$
\item $f'''(x)=\cos x$
\item $f'''(x)=8e^{-2x}-\sin x$
\item A car is being driven at a rate of $40$ mph when the brakes are applied. The car decelerates at a constant rate of $10$ ft/sec$^{2}$. How long before the car stops?
\item In the preceding problem, calculate how far the car travels in the time it takes to stop.
\item You are merging onto the freeway, accelerating from rest at a constant rate of $12$ ft/sec$^{2}$. How long does it take you to reach merging speed at $60$ mph?
\item Based on the previous problem, how far does the car travel to reach merging speed?
\item A car company wants to ensure its newest model can stop in $8$ sec when traveling at $75$ mph. If we assume constant deceleration, find the value of deceleration that accomplishes this.
\item A car company wants to ensure its newest model can stop in less than $450$ ft when traveling at $60$ mph. If we assume constant deceleration, find the value of deceleration that accomplishes this.
\end{enumerate}

For the following exercises, find the antiderivative of the function, assuming $F(0)=0$.

\begin{enumerate}[resume]
\item \techrequired $f(x)=x^{2}+2$
\item \techrequired $f(x)=4x-\sqrt{x}$
\item \techrequired $f(x)=\sin x+2x$
\item \techrequired $f(x)=e^{x}$
\item \techrequired $f(x)=\frac{1}{{(x+1)}^{2}}$
\item \techrequired $f(x)=e^{-2x}+3x^{2}$
\end{enumerate}

For the following exercises, determine whether the statement is true or false. Either prove it is true or find a counterexample if it is false.

\begin{enumerate}[resume]
\item If $f(x)$ is the antiderivative of $v(x)$, then $2f(x)$ is the antiderivative of $2v(x)$.
\item If $f(x)$ is the antiderivative of $v(x)$, then $f(2x)$ is the antiderivative of $v(2x)$.
\item If $f(x)$ is the antiderivative of $v(x)$, then $f(x)+1$ is the antiderivative of $v(x)+1$.
\item If $f(x)$ is the antiderivative of $v(x)$, then ${(f(x))}^{2}$ is the antiderivative of ${(v(x))}^{2}$.
\end{enumerate}

\end{document}
